% Generated draft: requires mathematical and source review.
\documentclass[11pt]{article}
\usepackage[margin=25mm]{geometry}
\usepackage{fontspec}
\setmainfont{Latin Modern Roman}
\usepackage{amsmath,amssymb,mathtools,unicode-math}
\setmathfont{Latin Modern Math}
\usepackage{enumitem,needspace,xurl,hyperref}
\setlength{\parindent}{0pt}
\setlength{\parskip}{5pt}
\setlength{\emergencystretch}{3em}
\newcommand{\N}{\mathbb{N}}
\newcommand{\Z}{\mathbb{Z}}
\newcommand{\Q}{\mathbb{Q}}
\newcommand{\R}{\mathbb{R}}
\newcommand{\C}{\mathbb{C}}
\newcommand{\F}{\mathbb{F}}
\newcommand{\A}{\mathbb{A}}
\newcommand{\PP}{\mathbb{P}}
\newcommand{\E}{\mathbb{E}}
\newcommand{\eps}{\varepsilon}
\newcommand{\dd}{\,\mathrm{d}}
\begin{document}
\section*{UnsolvedMath ID: 20000820}
\textbf{Source problem code:} catalogue entry
\subsection{Definitions and mathematical statement}
The source asks which numerical invariants occur for locally Cohen--Macaulay surfaces in $\PP^4$, naming degree and sectional genus. Work geometrically over an algebraically closed field, since the source specifies no field. Let $S\subset\PP^4$ be pure two-dimensional and locally Cohen--Macaulay, let $d=\deg S$, let $\pi=p_a(S\cap H)$ for a general regular hyperplane, and let $\chi=\chi(\mathcal O_S)$. Its Hilbert polynomial is
\[
 P_S(m)=\frac d2m^2+\left(\frac d2+1-\pi\right)m+\chi.
\]
The source does not require smoothness, integrality, nondegeneracy or arithmetic Cohen--Macaulayness.

The aid reports a forbidden strip, obtained from the regular hyperplane curve and Hartshorne's locally Cohen--Macaulay space-curve theorem: for $d\geq2$,
\[
 \pi=\binom{d-1}{2}\quad\hbox{or}\quad
 \pi\leq\binom{d-2}{2}.
\]
Degree one gives the plane case with sectional genus zero. This necessary restriction does not classify all possible triples.

For the arithmetically Cohen--Macaulay (ACM), codimension-two subgeography, write a Hilbert--Burch resolution
\[
 0\longrightarrow\bigoplus_j\mathcal O(-b_j)
 \longrightarrow\bigoplus_i\mathcal O(-a_i)
 \longrightarrow\mathcal I_S\longrightarrow0.
\]
With $M_r=\sum_j b_j^r-\sum_i a_i^r$, and the rank/first-moment cancellations of this resolution, the aid's moment formulas are
\[
 d=\frac{M_2}{2},\qquad
 \pi=1+\frac{M_3-6M_2}{6},\qquad
 \chi=\frac{M_4-10M_3+35M_2}{24}.
\]
These formulas assume an actual codimension-two ACM resolution; arbitrary degree lists need not define a surface.

For every pair of positive integers $(t,e)$, the report gives a realizable family from maximal minors of a general $t\times(t+1)$ matrix of degree-$e$ forms, of expected codimension two. The exact invariants are
\[
 d=\frac{t(t+1)e^2}{2},\qquad
 \pi=1+\frac{t(t+1)e^2((2t+1)e-6)}6,
\]
\[
 \chi=\frac{t(t+1)e^2\bigl((3t^2+3t+1)e^2-10(2t+1)e+35\bigr)}{24}.
\]
These are ACM, hence locally Cohen--Macaulay, examples; no assertion that every locally Cohen--Macaulay surface is ACM is made.

The aid further reports generation of new triples by sufficiently positive proper complete-intersection linkage. In the standard residual-link convention, let $S,S'$ be properly linked in a complete intersection $Y$ of degrees $(a,b)$, set $c=a+b-5$, and use the required local Cohen--Macaulay hypotheses. The transformation can be expressed by
\[
 P_{S'}(m)=P_Y(m)-P_S(c-m),
\]
hence
\[
 d'=ab-d,\qquad
 \pi'=\pi+\frac{(ab-2d)(a+b-4)}2,
\]
\[
 \chi'=\chi(\mathcal O_Y)-\chi-\frac d2c^2-c\left(\frac d2+1-\pi\right),
\]
where
\[
 \chi(\mathcal O_Y)=\frac{ab\bigl(2a^2+3ab+2b^2-15a-15b+35\bigr)}{12}.
\]
A proper residual link must actually exist; these identities are not an unrestricted recipe for declaring arbitrary numerical triples realizable. The forbidden strip, determinantal chart and linkage operations leave the full geography problem unresolved in the supplied assessment.
\subsection{Short English statement}
The aid gives a necessary gap in possible sectional genera and an explicit two-parameter family of realizable degree, genus and Euler-characteristic triples. Hilbert--Burch resolutions and proper linkage provide further computations. These cover parts of the geography, rather than all locally Cohen--Macaulay surfaces in projective four-space.
\subsection{Sources}
\begin{itemize}
\item Local UnsolvedMath registry, permanent ID $20000820$, source code \texttt{catalogue entry}, supplied complete source record and research summary.
\item American Institute of Mathematics, \emph{Components of Hilbert Schemes}, Problem 19, supplied URL \url{https://aimath.org/WWN/hilbertschemes/hilbertschemes.pdf}; canonical local record \texttt{aim-arithmetic-geometry-notes.json}, \texttt{notes[65]}. No PDF was retrieved.
\item The summary names Hartshorne's locally Cohen--Macaulay space-curve theorem, Hilbert--Burch resolutions and complete-intersection linkage but supplies no full bibliographic citations. The original literature field is empty.
\item Research attempt 1: partial result; problem open at the recorded run, registry status partially solved; supplied rigor label ``proved'', novelty confidence ``low''. The machine-generated aid is not peer reviewed and requires independent verification before citation.
\end{itemize}
\end{document}
